\documentclass[14pt, a4paper]{extarticle}
\usepackage[T2A]{fontenc}
\usepackage[utf8]{inputenc}
\usepackage[english, russian]{babel}
\usepackage{tempora}
\usepackage{longtable}
\usepackage{array}
\usepackage{tabularx}
\usepackage{ltablex}
\usepackage{ragged2e}
\usepackage{graphicx}
\usepackage{calc}
\usepackage[left=30mm, right=15mm, top=20mm, bottom=20mm, footskip=12mm]{geometry}

\setlength{\parindent}{0pt}
\setlength{\parskip}{6pt plus 2pt minus 1pt}
\keepXColumns

\newcolumntype{P}[1]{>{\RaggedRight\arraybackslash}p{#1 - 2\tabcolsep - 2\arrayrulewidth}}
\newcolumntype{C}[1]{>{\centering\arraybackslash}p{#1 - 2\tabcolsep - 2\arrayrulewidth}}

\begin{document}
\pagestyle{empty}

\begin{center}
    {\small МИНИСТЕРСТВО НАУКИ И ВЫСШЕГО ОБРАЗОВАНИЯ РОССИЙСКОЙ ФЕДЕРАЦИИ}

    \resizebox{\textwidth}{!}{\fontsize{8pt}{10pt}\selectfont ФЕДЕРАЛЬНОЕ
        ГОСУДАРСТВЕННОЕ АВТОНОМНОЕ ОБРАЗОВАТЕЛЬНОЕ УЧРЕЖДЕНИЕ ВЫСШЕГО ОБРАЗОВАНИЯ}

    {\small \textbf{«МОСКОВСКИЙ ПОЛИТЕХНИЧЕСКИЙ УНИВЕРСИТЕТ»}} \\
    {\small \textbf{(МОСКОВСКИЙ ПОЛИТЕХ)}}
\end{center}

\begin{tabularx}{\textwidth}{@{} p{0.53\textwidth} p{0.43\textwidth} @{} }
     &
    \raggedright
    УТВЕРЖДАЮ                          \\
    заведующая кафедрой                \\
    «Инфокогнитивные технологии»       \\[0.3cm]

    \rule{3cm}{0.4pt} / Е.~А.~Пухова / \\[0.1cm]

    «\rule{0.8cm}{0.4pt}» \rule{3cm}{0.4pt} 2026~г.
\end{tabularx}

\begin{center}
    {\large \textbf{ЗАДАНИЕ}}

    \textbf{на выполнение курсовой работы (проекта)}

    Николаеву Алексею Владимировичу
\end{center}

обучающемуся группы 241-3211,

направления подготовки 09.03.01 «Информатика и вычислительная техника»

по дисциплине «Разработка веб-приложений»

на тему «Разработка веб-приложения для электронного архива учебных материалов с ролевым доступом»

1. Исходные данные к работе (проекту): информационные ресурсы в сети интернет, научные публикации в открытой печати, документация по средствам разработки веб-приложений.

2. Содержание задания по курсовой работе (проекту) - перечень вопросов, подлежащих разработке:

{\renewcommand{\arraystretch}{1.18}
\small
\begin{longtable}{|P{0.50\textwidth}|C{0.11\textwidth}|C{0.15\textwidth}|P{0.24\textwidth}|}
\hline
\textbf{Разрабатываемый вопрос} & \textbf{Объем, \%} & \textbf{Срок} & \textbf{Примечание} \\
\hline
\endhead
\hline
\endfoot

\multicolumn{4}{|P{\textwidth}|}{\textbf{Раздел 1. Анализ предметной области}} \\ \hline
Задача 1.1. Анализ существующих веб-приложений для хранения, систематизации и публикации учебных материалов & 10 & 06.04.2026 & \\ \hline
Задача 1.2. Анализ программных средств разработки и выбор стека технологий & 7 & 12.04.2026 & \\ \hline
Задача 1.3. Формулировка цели, задач и требований к системе & 8 & 18.04.2026 & \\ \hline

\multicolumn{4}{|P{\textwidth}|}{\textbf{Раздел 2. Проектирование веб-приложения}} \\ \hline
Задача 2.1. Определение ролей пользователей и прав доступа (администратор, преподаватель, студент) & 6 & 25.04.2026 & обязательный пункт \\ \hline
Задача 2.2. Проектирование функциональности приложения и сценариев использования & 7 & 03.05.2026 & \\ \hline
Задача 2.3. Проектирование модели данных для сущностей «Пользователь», «Категория», «Материал» и «Файл» & 9 & 10.05.2026 & не менее 3 сущностей \\ \hline
Задача 2.4. Проектирование структуры интерфейса и макетов основных страниц & 5 & 16.05.2026 & \\ \hline

\multicolumn{4}{|P{\textwidth}|}{\textbf{Раздел 3. Разработка веб-приложения}} \\ \hline
Задача 3.1. Разработка базовой структуры приложения и шаблонов страниц & 8 & 20.05.2026 & \\ \hline
Задача 3.2. Реализация аутентификации и авторизации пользователей & 7 & 24.05.2026 & обязательный пункт \\ \hline
Задача 3.3. Реализация системы проверки прав доступа в соответствии с ролью пользователя & 7 & 28.05.2026 & обязательный пункт \\ \hline
Задача 3.4. Реализация CRUD-интерфейса для пользователей, категорий и материалов & 10 & 01.06.2026 & обязательный пункт \\ \hline
Задача 3.5. Реализация загрузки, скачивания и удаления файлов учебных материалов & 8 & 04.06.2026 & обязательный пункт \\ \hline
Задача 3.6. Реализация агрегирования данных и формирования статистики по материалам & 6 & 06.06.2026 & количество файлов, материалов по категориям и авторам \\ \hline

\multicolumn{4}{|P{\textwidth}|}{\textbf{Раздел 4. Оформление итогов работы}} \\ \hline
Задача 4.1. Тестирование основных пользовательских сценариев и проверка прав доступа & 4 & 09.06.2026 & \\ \hline
Задача 4.2. Подготовка Git-репозитория и итоговой версии исходного кода проекта & 3 & 11.06.2026 & \\ \hline
Задача 4.3. Подготовка демонстрационной версии приложения и материалов к защите & 3 & 12.06.2026 & \\ \hline
Задача 4.4. Оформление пояснительной записки в соответствии с требованиями & 3 & 13.06.2026 & \\ \hline
\end{longtable}
}

\clearpage

Руководитель курсовой работы (проекта): \\ старший преподаватель кафедры «Инфокогнитивные технологии»

{\renewcommand{\arraystretch}{2}
\noindent
\begin{tabularx}{\textwidth}{@{} P{0.45\textwidth} P{0.30\textwidth} P{0.25\textwidth} @{} }
«\rule{0.8cm}{0.4pt}» \quad \rule{2cm}{0.4pt} \quad 2026~г. &
\centering \rule{3.5cm}{0.4pt} \par \footnotesize (подпись) &
\hfill А.~С.~Кружалов \\
\end{tabularx}

\vspace{1.2cm}

\begin{tabularx}{\textwidth}{@{} P{0.52\textwidth} P{0.48\textwidth} @{} }
Дата выдачи задания & \hfill 02.04.2026~г. \\

Дата сдачи выполненной работы & \hfill «\rule{0.8cm}{0.4pt}» \quad \rule{2cm}{0.4pt} \quad 2026~г. \\
\end{tabularx}
}

\vspace{1.6cm}

Задание принял к исполнению

{\renewcommand{\arraystretch}{1.5}
\begin{tabularx}{\textwidth}{@{} P{0.45\textwidth} P{0.30\textwidth} P{0.25\textwidth} @{} }
02.04.2026~г. &
\centering \rule{3.5cm}{0.4pt} \par \footnotesize (подпись) &
\hfill Николаев А.~В. \\
\end{tabularx}
}

\end{document}